\chapter[Ví dụ minh họa]{Ví dụ minh họa}
\label{cp:example}

\section{Ví dụ 1: Cơ sở}

\begin{figure}[htbp]
    \centering
    \renewcommand{\arraystretch}{1.2} % Tăng nhẹ khoảng cách các dòng trong bảng

    % ================= HÀNG 1 =================
    \begin{minipage}[t]{0.22\textwidth}
        \centering
        CSDL\\[4pt]
        \begin{tabular}{|c|l|}
        \hline
        TID & Items \\ \hline
        101 & 1 2 3 4 \\
        102 & 1 2 \\
        103 & 1 4 5 \\
        104 & 1 2 4 5 \\
        105 & 5 \\ \hline
        \end{tabular}
    \end{minipage}
    \hfill
    \begin{minipage}[t]{0.45\textwidth}
        \centering
        $\overline{C}_1$\\[4pt]
        \begin{tabular}{|c|l|}
        \hline
        TID & Tập Itemset \\ \hline
        101 & \{\{1\}, \{2\}, \{3\}, \{4\}\} \\
        102 & \{\{1\}, \{2\}\} \\
        103 & \{\{1\}, \{4\}, \{5\}\} \\
        104 & \{\{1\}, \{2\}, \{4\}, \{5\}\} \\
        105 & \{5\} \\ \hline
        \end{tabular}
    \end{minipage}
    \hfill
    \begin{minipage}[t]{0.25\textwidth}
        \centering
        $L_1$\\[4pt]
        \begin{tabular}{|c|c|}
        \hline
        Itemset & Support \\ \hline
        \{1\} & 4 \\
        \{2\} & 3 \\
        \{4\} & 3 \\
        \{5\} & 3 \\ \hline
        \end{tabular}
    \end{minipage}

    \vspace{0.8cm} % Khoảng cách giữa hàng 1 và hàng 2

    % ================= HÀNG 2 =================
    \begin{minipage}[t]{0.22\textwidth}
        \centering
        $C_2$\\[4pt]
        \begin{tabular}{|c|}
        \hline
        Itemset \\ \hline
        \{1 2\} \\
        \{1 4\} \\
        \{1 5\} \\
        \{2 4\} \\
        \{2 5\} \\
        \{4 5\} \\ \hline
        \end{tabular}
    \end{minipage}
    \hfill
    \begin{minipage}[t]{0.45\textwidth}
        \centering
        $\overline{C}_2$\\[4pt]
        \begin{tabular}{|c|l|}
        \hline
        TID & Tập Itemset \\ \hline
        101 & \{\{1 2\}, \{1 4\}, \{2 4\}\} \\
        102 & \{\{1 2\}\} \\
        103 & \{\{1 4\}, \{1 5\}, \{4 5\}\} \\
        104 & \begin{tabular}{@{}l@{}}\{\{1 2\}, \{1 4\}, \{1 5\}, \\ \{2 4\}, \{2 5\}, \{4 5\}\}\end{tabular} \\ \hline
        \end{tabular}
    \end{minipage}
    \hfill
    \begin{minipage}[t]{0.25\textwidth}
        \centering
        $L_2$\\[4pt]
        \begin{tabular}{|c|c|}
        \hline
        Itemset & Support \\ \hline
        \{1 2\} & 3 \\
        \{1 4\} & 3 \\
        \{1 5\} & 2 \\
        \{2 4\} & 2 \\
        \{4 5\} & 2 \\ \hline
        \end{tabular}
    \end{minipage}

    \vspace{0.8cm} % Khoảng cách giữa hàng 2 và hàng 3

    % ================= HÀNG 3 =================
    \begin{minipage}[t]{0.22\textwidth}
        \centering
        $C_3$\\[4pt]
        \begin{tabular}{|c|}
        \hline
        Itemset \\ \hline
        \{1 2 4\} \\
        \{1 2 5\} \\
        \{1 4 5\} \\
        \{2 4 5\} \\ \hline
        \end{tabular}
    \end{minipage}
    \hfill
    \begin{minipage}[t]{0.45\textwidth}
        \centering
        $\overline{C}_3$\\[4pt]
        \begin{tabular}{|c|l|}
        \hline
        TID & Tập Itemset \\ \hline
        101 & \{\{1 2 4\}\} \\
        103 & \{\{1 4 5\}\} \\
        104 & \begin{tabular}{@{}l@{}}\{\{1 2 4\}, \{1 2 5\}, \\ \{1 4 5\}, \{2 4 5\}\}\end{tabular} \\ \hline
        \end{tabular}
    \end{minipage}
    \hfill
    \begin{minipage}[t]{0.25\textwidth}
        \centering
        $L_3$\\[4pt]
        \begin{tabular}{|c|c|}
        \hline
        Itemset & Support \\ \hline
        \{1 2 4\} & 2 \\
        \{1 4 5\} & 2 \\ \hline
        \end{tabular}
    \end{minipage}

    \caption{Ví dụ 1 --- Các bước thực thi AprioriTID trên CSDL 5 giao dịch với $\text{minsup}=2$.}
    \label{fig:example-1}
\end{figure}

\textbf{Dừng tại $k=4$:} Áp dụng \textit{apriori-gen} trên $L_3 = \{\{1,2,4\},\{1,4,5\}\}$, bước \textit{join} yêu cầu hai itemset phải trùng $k-2 = 2$ item đầu. Ở đây $\{1,2,4\}[1{:}2] = \{1,2\}$ khác $\{1,4,5\}[1{:}2] = \{1,4\}$, nên không sinh được ứng viên nào, tức $C_4 = \emptyset$. Thuật toán kết thúc.

\textbf{Tập kết quả cuối cùng} (với support tương ứng):
\begin{center}
\begin{tabular}{|l|c||l|c|}
\hline
Itemset & Support & Itemset & Support \\ \hline
\{1\}     & 4 & \{1,5\}   & 2 \\
\{2\}     & 3 & \{2,4\}   & 2 \\
\{4\}     & 3 & \{4,5\}   & 2 \\
\{5\}     & 3 & \{1,2,4\} & 2 \\
\{1,2\}   & 3 & \{1,4,5\} & 2 \\
\{1,4\}   & 3 & & \\ \hline
\end{tabular}
\end{center}

Tổng cộng thu được \textbf{11 tập phổ biến}.

\textbf{Kiểm tra chéo (cross-check) bằng liệt kê thủ công.} Duyệt toàn bộ CSDL gốc:
\begin{itemize}[nosep,topsep=2pt]
    \item \textbf{1-itemset:} \{1\} xuất hiện trong T101, T102, T103, T104 $\Rightarrow$ sup=4. \{2\} trong T101, T102, T104 $\Rightarrow$ sup=3. \{3\} chỉ trong T101 $\Rightarrow$ sup=1 (\textit{loại}). \{4\} trong T101, T103, T104 $\Rightarrow$ sup=3. \{5\} trong T103, T104, T105 $\Rightarrow$ sup=3.
    \item \textbf{2-itemset phổ biến:} \{1,2\} (T101,T102,T104)=3; \{1,4\} (T101,T103,T104)=3; \{1,5\} (T103,T104)=2; \{2,4\} (T101,T104)=2; \{4,5\} (T103,T104)=2. Cặp \{2,5\} chỉ xuất hiện trong T104 nên bị loại.
    \item \textbf{3-itemset phổ biến:} \{1,2,4\} (T101,T104)=2; \{1,4,5\} (T103,T104)=2. Các tổ hợp khác hoặc chứa item không phổ biến, hoặc có support $<2$.
    \item \textbf{4-itemset:} \{1,2,4,5\} chỉ xuất hiện trong T104 $\Rightarrow$ sup=1 (\textit{loại}).
\end{itemize}
Kết quả liệt kê thủ công đúng \textbf{11 tập phổ biến}, khớp hoàn toàn với đầu ra của thuật toán ở \autoref{fig:example-1}.

\section{Ví dụ 2: Tình huống đặc biệt}

Điểm yếu nổi tiếng nhất của AprioriTID là ở giai đoạn $k=2$, tập $\overline{C}_2$ có thể \textbf{lớn hơn rất nhiều} so với CSDL gốc $\mathcal{D}$. Điều này xảy ra khi dữ liệu dày đặc (dense): mỗi giao dịch chứa nhiều item, sinh ra số lượng cặp ứng viên rất lớn.

Ta thiết kế CSDL $\mathcal{D}'$ gồm 4 giao dịch, mỗi giao dịch chứa 6 item từ tập 7 item, với $\text{minsup} = 2$:

\begin{figure}[htbp]
    \centering
    \renewcommand{\arraystretch}{1.2}

    % ================= HÀNG 1: DB, C̄_1, L_1 =================
    \begin{minipage}[t]{0.22\textwidth}
        \centering
        Database $\mathcal{D}'$\\[4pt]
        \begin{tabular}{|c|l|}
        \hline
        TID & Items \\ \hline
        101 & 1 2 3 4 5 6 \\
        102 & 1 2 3 4 5 7 \\
        103 & 1 2 3 4 6 7 \\
        104 & 1 2 3 5 6 7 \\ \hline
        \end{tabular}
    \end{minipage}
    \hfill
    \begin{minipage}[t]{0.45\textwidth}
        \centering
        $\overline{C}_1$\\[4pt]
        \begin{tabular}{|c|l|}
        \hline
        TID & Tập Itemset \\ \hline
        101 & \{\{1\},\{2\},\{3\},\{4\},\{5\},\{6\}\} \\
        102 & \{\{1\},\{2\},\{3\},\{4\},\{5\},\{7\}\} \\
        103 & \{\{1\},\{2\},\{3\},\{4\},\{6\},\{7\}\} \\
        104 & \{\{1\},\{2\},\{3\},\{5\},\{6\},\{7\}\} \\ \hline
        \end{tabular}
    \end{minipage}
    \hfill
    \begin{minipage}[t]{0.25\textwidth}
        \centering
        $L_1$\\[4pt]
        \begin{tabular}{|c|c|}
        \hline
        Itemset & Support \\ \hline
        \{1\} & 4 \\
        \{2\} & 4 \\
        \{3\} & 4 \\
        \{4\} & 3 \\
        \{5\} & 3 \\
        \{6\} & 3 \\
        \{7\} & 3 \\ \hline
        \end{tabular}
    \end{minipage}

    \caption{Ví dụ 2 --- Bước $k=1$: Tất cả 7 item đều phổ biến.}
    \label{fig:example-2-k1}
\end{figure}

\textbf{Phân tích kích thước $\overline{C}_1$:} CSDL gốc $\mathcal{D}'$ có $4 \times 6 = 24$ mục (item entries). Tập $\overline{C}_1$ cũng chứa đúng $4 \times 6 = 24$ mục (1-itemset entries) vì tất cả item đều phổ biến, không có item nào bị loại.

\textbf{Bước $k=2$:} Do $|L_1| = 7$, hàm apriori-gen sinh ra $C_2 = \binom{7}{2} = 21$ ứng viên 2-itemset. Mỗi giao dịch chứa 6 item, tạo ra $\binom{6}{2} = 15$ cặp con hợp lệ. Tổng số mục trong $\overline{C}_2$:

$$|\overline{C}_2|_{\text{entries}} = 4 \times 15 = 60 \quad \text{vs} \quad |\mathcal{D}'|_{\text{entries}} = 24$$

\textbf{$\overline{C}_2$ lớn gấp 2.5 lần CSDL gốc!} Đây chính là điểm yếu bộ nhớ đặc trưng của AprioriTID trên dữ liệu dày đặc.

\begin{figure}[htbp]
    \centering
    \renewcommand{\arraystretch}{1.2}

    \begin{minipage}[t]{0.30\textwidth}
        \centering
        $C_2$ (21 ứng viên)\\[4pt]
        \begin{tabular}{|c|}
        \hline
        Itemset \\ \hline
        \{1 2\}, \{1 3\}, \{1 4\} \\
        \{1 5\}, \{1 6\}, \{1 7\} \\
        \{2 3\}, \{2 4\}, \{2 5\} \\
        \{2 6\}, \{2 7\}, \{3 4\} \\
        \{3 5\}, \{3 6\}, \{3 7\} \\
        \{4 5\}, \{4 6\}, \{4 7\} \\
        \{5 6\}, \{5 7\}, \{6 7\} \\ \hline
        \end{tabular}
    \end{minipage}
    \hfill
    \begin{minipage}[t]{0.65\textwidth}
        \centering
        $\overline{C}_2$ (\textbf{60 entries} --- lớn hơn $\mathcal{D}'$)\\[4pt]
        \begin{tabular}{|c|l|}
        \hline
        TID & Tập Itemset (15 cặp mỗi giao dịch) \\ \hline
        101 & \begin{tabular}{@{}l@{}}\{\{1 2\},\{1 3\},\{1 4\},\{1 5\},\{1 6\},\{2 3\},\{2 4\},\\ \{2 5\},\{2 6\},\{3 4\},\{3 5\},\{3 6\},\{4 5\},\{4 6\},\{5 6\}\}\end{tabular} \\
        102 & \begin{tabular}{@{}l@{}}\{\{1 2\},\{1 3\},\{1 4\},\{1 5\},\{1 7\},\{2 3\},\{2 4\},\\ \{2 5\},\{2 7\},\{3 4\},\{3 5\},\{3 7\},\{4 5\},\{4 7\},\{5 7\}\}\end{tabular} \\
        103 & \begin{tabular}{@{}l@{}}\{\{1 2\},\{1 3\},\{1 4\},\{1 6\},\{1 7\},\{2 3\},\{2 4\},\\ \{2 6\},\{2 7\},\{3 4\},\{3 6\},\{3 7\},\{4 6\},\{4 7\},\{6 7\}\}\end{tabular} \\
        104 & \begin{tabular}{@{}l@{}}\{\{1 2\},\{1 3\},\{1 5\},\{1 6\},\{1 7\},\{2 3\},\{2 5\},\\ \{2 6\},\{2 7\},\{3 5\},\{3 6\},\{3 7\},\{5 6\},\{5 7\},\{6 7\}\}\end{tabular} \\ \hline
        \end{tabular}
    \end{minipage}

    \caption{Ví dụ 2 --- Bước $k=2$: $\overline{C}_2$ có 60 entries, gấp 2.5 lần CSDL gốc (24 entries).}
    \label{fig:example-2-k2}
\end{figure}

Đếm support từ $\overline{C}_2$, ta được $L_2$ gồm tất cả 21 cặp đều phổ biến (support $\geq 2$):

\begin{itemize}
    \item Support 4: \{1 2\}, \{1 3\}, \{2 3\}
    \item Support 3: \{1 4\}, \{1 5\}, \{1 6\}, \{1 7\}, \{2 4\}, \{2 5\}, \{2 6\}, \{2 7\}, \{3 4\}, \{3 5\}, \{3 6\}, \{3 7\}
    \item Support 2: \{4 5\}, \{4 6\}, \{4 7\}, \{5 6\}, \{5 7\}, \{6 7\}
\end{itemize}

Các bước tiếp theo ($k=3, 4, 5$) tương tự. Tại $k=3$: $|C_3| = \binom{7}{3} = 35$ ứng viên, $|L_3| = 31$ itemset phổ biến. Tại $k=4$: $|L_4| = 22$. Tại $k=5$: $|L_5| = 6$. Tổng cộng thu được $7 + 21 + 31 + 22 + 6 = \textbf{87 tập phổ biến}$.

\textbf{Tại sao đây là tình huống đặc biệt quan trọng?}

\begin{itemize}
    \item Với dữ liệu dày đặc (dense), mỗi giao dịch có $w$ item sẽ sinh ra $\binom{w}{2}$ cặp trong $\overline{C}_2$. Khi $w$ lớn, $\overline{C}_2$ tăng \textbf{bậc hai} theo $w$, trong khi $\mathcal{D}$ chỉ tăng tuyến tính.
    \item Đây là lý do chính khiến AprioriTID có thể tiêu tốn nhiều bộ nhớ hơn Apriori chuẩn (vốn quét lại $\mathcal{D}$ gốc), đặc biệt khi minsup thấp và dữ liệu dày đặc.
    \item Tuy nhiên, ở các bước $k$ lớn hơn, $\overline{C}_k$ giảm nhanh do cơ chế transaction reduction --- các giao dịch không chứa ứng viên bị loại bỏ. Lợi thế của AprioriTID thể hiện rõ ở các bước $k$ cao, khi $\overline{C}_k \ll |\mathcal{D}|$.
\end{itemize}

\textbf{Kiểm tra chéo:} Liệt kê thủ công --- mỗi giao dịch thiếu đúng 1 item (T101 thiếu 7, T102 thiếu 6, T103 thiếu 5, T104 thiếu 4). Một tập con $S$ có support bằng $4 - |\{$item thiếu nằm trong $S\}|$. Bất kỳ tập con nào cũng xuất hiện trong ít nhất 2 giao dịch (vì chỉ có 4 giao dịch, mỗi giao dịch thiếu đúng 1 item khác nhau từ \{4,5,6,7\}). Riêng các tập con chỉ chứa \{1,2,3\} luôn có support = 4. Kết quả 87 tập phổ biến khớp với kết quả chạy thuật toán.